% Full draft: text + lean table set. Compile with pdflatex -> bibtex -> pdflatex x2 (TinyTeX).
% Text: main_text_part1.tex (Sections 1-3), main_text_part2.tex (Sections 4-5). Tables: tables_tax.tex (generated by 111).
\documentclass[11pt]{article}
\usepackage[margin=2.5cm]{geometry}
\usepackage{amsmath,booktabs,threeparttable,rotating,makecell,caption,graphicx,adjustbox}
\usepackage[authoryear,round]{natbib}
\usepackage[hidelinks]{hyperref}
\usepackage{setspace}
\onehalfspacing
\captionsetup{labelfont = bf, textfont = it, justification = centering}
\title{Taxing the periphery:\\ air passenger taxes, network position and aviation CO$_2$ in Europe}
\author{Sunbin Yoo\thanks{Department of Energy, Sungkyunkwan University.} \and [coauthors]}
\date{Draft, October 2026}
\begin{document}
\maketitle
\begin{abstract}
\noindent Nine European countries introduced or raised per-passenger air ticket taxes between 1998 and 2018.
Using a monthly panel of more than 6,000 airports (1996--2024), flight-stage CO$_2$ built from the same
schedules, and a network-centrality index of each airport's position in the world air network, we ask
whether these taxes shrink a country's air network or concentrate it. Following the local comparison
design of Li, Liu, Purevjav and Yang (2019), we compare airports in the taxing country with airports in
European countries without a tax event around each verified announcement, keep foreign airports near the
border as a separate exposed group, and test parallel trends with event studies, synthetic
difference-in-differences and sensitivity analysis. Taxed airports lose 10 percent of seats and CO$_2$ in
the first tax year with equal weights but 4 and 2 percent with seat weights; the smallest third of airports
lose 24 percent of seats, 14 percent of destinations and a third of their eigenvector centrality, while
hubs lose nothing and upgauge. National seat totals and national connectivity are unchanged for six of nine
events. Emissions fall only through route exit; emissions per seat-kilometre do not fall. Leakage across
borders appears only in the Dutch episode. Revenue per tonne avoided is about EUR 1,000, and the abatement
is paid in the connectivity of peripheral airports.\\[4pt]
\textit{Keywords}: aviation tax; carbon pricing; air connectivity; network centrality; difference-in-differences; leakage\\
\textit{JEL}: H23, Q54, R41, L93
\end{abstract}
% ---------------------------------------------------------------------------------------------
% Main text, part 1 (Sections 1-3). Draft written 2026-10-03 from the lean table set (tables_tax.tex)
% and the result CSVs in GACI_FuelShock/. Structure follows Li, Liu, Purevjav and Yang (2019, JEEM).
% Numbers in the text are taken from tables_tax.tex / _res_tax_*.csv; re-check after any re-run of 92-111.
% ---------------------------------------------------------------------------------------------

\section{Introduction}

Air passenger taxes are the most widely used national instrument for pricing the climate cost of flying.
Nine European countries introduced or raised a per-passenger ticket tax between 1998 and 2018, several more
have followed since, and the European Commission's Energy Taxation Directive proposal and the 2027 Dutch
distance-band reform keep the instrument on the agenda. The tax is simple to collect and raises
substantial revenue, but what it does to flying is less clear. A per-passenger levy can reduce emissions
by cutting the number of flights; it can also redistribute flying within the network, because every
European ticket tax exempts transfer passengers and is a larger share of the fare on short, cheap
routes than on long ones. Whether a ticket tax shrinks a country's air network or merely concentrates it
is an empirical question, and the answer determines both the carbon abatement the tax delivers and who
pays for it.

This paper answers that question with the first evaluation of ticket taxes that treats the air network
itself as the outcome. We assemble a monthly panel of scheduled seats and flights for more than 6,000
airports from 1996 to 2024, build flight-stage-resolved CO$_2$ emissions for every departing flight from the
same schedules, and attach to each airport-year a network-centrality summary of its position in the
world air network, the Global Air Connectivity Index (GACI) of \citet{cheung2020}. We then evaluate the
nine European ticket-tax events listed in Table~\ref{tab:events}: the extension of the Danish passenger
tax to domestic flights (1998), the Maltese departure-tax increase (2005), the doubling of UK Air
Passenger Duty (2007), the Dutch ticket tax (2008, repealed in 2009), the Irish Air Travel Tax (2009),
the German and Austrian taxes (2011), and the Norwegian (2016) and Swedish (2018) taxes.

Our empirical framework follows \citet{li2019subway}, who evaluate Beijing's subway expansion with two
complementary designs: a difference-in-differences comparison of locations near and far from each
opening, with a buffer zone to keep misclassified locations out of the control group, and a continuous
network-density measure that captures system-wide effects. We transpose both to the air network. The
difference-in-differences design compares airports in the taxing country with airports in European
countries that have no tax event in the same window, around each tax's announcement. Airports in
untaxed countries that lie within 50 to 300 kilometres of the taxed border enter as a separate exposed
group rather than as controls, because leakage across the border would otherwise contaminate the
comparison, and airports within 50 kilometres are dropped. Event time runs from the announcement, with
the months between announcement and implementation removed, so that anticipation cannot enter either
the pre- or the post-period. The continuous design uses GACI as the network-density analogue: we
estimate how an airport's connectivity moves when its country is taxed, how emissions respond to
connectivity, and how much of the emissions effect runs through the network.

Three results emerge. First, taxes reduce flying at the periphery and leave hubs intact. Across the
eight non-Dutch events, scheduled seats at taxed airports fall by 10 percent in the first tax year when
every airport counts equally, but by 4 percent when airports are weighted by seats. The gap is the
distribution of the effect: airports in the bottom size tercile of the taxing country lose 24 percent of
their seats, 14 percent of their destinations and 4 percentage points of their months with any service,
while large airports lose no seats and increase aircraft size by 2.5 percent. Second, the network
concentrates rather than shrinks. Taxed airports lose 1.6 to 2.3 percent of their connectivity index, 9
percent of their destinations, 19 percent of eigenvector centrality and 43 percent of betweenness
centrality, with the losses again concentrated in small and mid-sized airports; yet country-level
synthetic difference-in-differences estimates on national seat totals and on the seat-weighted national
connectivity index are indistinguishable from zero for every event except Ireland, Malta and Norway.
Third, emissions fall only where flights disappear. Departure CO$_2$ at taxed airports falls by 10
percent unweighted and 2 percent (not significant) seat-weighted; CO$_2$ per seat and per seat-kilometre
do not fall and rise slightly at the hubs that upgauge. An exact decomposition of the CO$_2$ change
attributes all of it to the route-exit margin, with no contribution from flight frequency on surviving
routes or from aircraft and distance. Cross-border leakage is absent in the pooled sample and
measurable only in the Dutch episode, where seats at foreign airports within 150 kilometres rose by 16
percent while the tax was in force.

The paper contributes to three literatures. The first evaluates ticket taxes directly. Synthetic-control,
panel and route-level studies of the German, Austrian, Dutch, Swedish and UK taxes find passenger or
flight reductions of 2 to 12 percent concentrated in low-cost and secondary airports and little or no
effect at hubs \citep{borbely2019,falk2019,bernardo2024flight,oesingmann2022,seetaram2014,gordijn2011}.
We confirm the size gradient on a common footing across nine events, add emissions and an exact
decomposition of how they fall, and show that the gradient is a reallocation of network position, not
just of volume. The second literature evaluates carbon pricing in aviation, chiefly the inclusion of
intra-EEA flights in the EU Emissions Trading System, which lowered emissions on regulated routes by
about 5 percent \citep{fageda2022pricing} and shifted intercontinental connecting passengers toward
hubs outside the system \citep{fageda2025hub}. Ticket taxes differ from allowance prices in two ways
that our results make precise: they are flat per passenger rather than proportional to fuel, so they do
not touch emissions intensity, and they exempt transfers, so they protect the hub function that carbon
pricing is found to erode. The third literature measures air networks with centrality indices
\citep{cheung2020} but has not used them as the outcome of a policy evaluation; the only causal designs
with centrality outcomes concern airline mergers. Our companion work estimates the emissions elasticity
of connectivity at the global level \citep{yoo2026co2} and its trade effects \citep{yoo2026trade}; here
we close the loop by estimating how a specific policy moves connectivity and how the connectivity
change maps into emissions.

The rest of the paper proceeds as in \citet{li2019subway}. Section~\ref{sec:background} describes the
tax events and the data. Section~\ref{sec:strategy} sets out the stacked difference-in-differences,
event-study, synthetic difference-in-differences and decomposition designs. Section~\ref{sec:results}
presents the results for seats, network position and emissions, the mechanism, the national totals,
leakage and the cost per tonne. Section~\ref{sec:conclusion} concludes.

\section{Background and data}
\label{sec:background}

\subsection{Ticket taxes in Europe}

Table~\ref{tab:events} lists the nine events. Every tax is levied per departing passenger on scheduled
commercial flights, is paid by the airline and passed into the fare, and exempts transfer passengers
who connect within a stated time at the taxed airport. Six taxes differentiate by destination distance
in two or three bands (Germany, Austria, Sweden, the Netherlands, Ireland and the United Kingdom); the
Norwegian, Danish and Maltese taxes are flat. Rates at introduction range from about EUR 2 on Irish
short-haul departures to EUR 45 on Dutch and German long-haul departures. Two features of the
institutional setting matter for the design. First, announcement and implementation are separated by
two to eighteen months, and in Germany the tax applied to tickets sold from September 2010 for flights
from January 2011, so demand could react before the effective date. We date events from the
announcement, verified in primary sources (budget speeches, draft bills and official gazettes), and
drop the months between announcement and implementation. Second, the taxes are national, so the
natural control group is airports in European countries without a tax event in the same window, and
the natural threat to that control group is leakage to foreign airports within driving distance of the
taxed country.

\subsection{Airports, schedules and emissions}

The unit of observation is the airport-month. Scheduled seats, flights and seat-kilometres for every
departing flight are taken from the OAG schedule database for January 1996 to June 2024, aggregated to
the airport and month and split into domestic and international departures. CO$_2$ emissions are
computed for each scheduled flight from aircraft-type fuel burn by flight stage (taxi-out, take-off,
climb-out and cruise) and the great-circle distance of the segment, then summed to the airport-month;
the inventory and its validation against EDGAR and national totals are described in
\citet{yoo2026co2}. Attributing cruise emissions to the departure airport avoids double counting and
follows the departure convention of national inventories. We use departure CO$_2$, CO$_2$ per seat and
CO$_2$ per seat-kilometre (emissions intensity), and the average stage length (seat-kilometres per
seat) as outcomes.

\subsection{Network position}

An airport's position in the world air network is measured by the Global Air Connectivity Index
\citep{cheung2020}, recomputed for every year from 1996 to 2024 on the full schedule network of
3,200 to 3,900 active airports. GACI is the first principal component of five standardised centrality
indicators: degree (number of destinations), closeness, eigenvector centrality (connections to
well-connected airports), flow betweenness (transfer flows routed through the airport, computed from an
electrical-circuit representation of the weighted network) and regional importance (intensity of
intra-regional links). The loadings are stable across years and the index is standardised so that
relative scores are comparable over time. We use the index and its degree, eigenvector and betweenness
components as annual outcomes. GACI is a relative measure of position within the year's network, so a
fall at taxed airports means they lost ground relative to other airports, not that the network shrank.

\subsection{Sample and descriptive patterns}

The estimation sample contains airports in 21 European countries for which tax histories were coded
from legal texts (Italy, whose municipal surcharge taxes transfer passengers and has no verified
dates, is excluded from both groups). For each event the treated group is every airport in the taxing
country with scheduled service in the year before the announcement; the control group is every airport
in a coded European country with no tax event within the window; airports in untaxed countries within
50 to 300 kilometres of the nearest treated airport form a border group, and those within 50
kilometres are dropped. Within each taxing country, airports are split into size terciles by seats in
the year before the announcement. The pooled sample has 112,876 airport-months in 33 countries.

A first look at the data follows \citet{li2019subway}. In the twelve months before and after the
effective date, log seats at taxed airports fall from 9.86 to 9.80 while seats at control airports rise
from 9.65 to 9.69. After partialling out airport-by-calendar-month fixed effects, the residual at taxed
airports moves from $+0.034$ to $-0.034$, at control airports from $-0.013$ to $+0.013$, and at border
airports from $-0.015$ to $+0.015$: a raw difference-in-differences of about 7 percent against
controls, with border airports moving with the controls rather than against them. The regressions below
sharpen this comparison.

\section{Empirical strategy}
\label{sec:strategy}

\subsection{Stacked difference-in-differences with a buffer}

For each event $e$ we build a stack of airport-months from 36 months before the announcement $A_e$ to
twelve months after the effective date $E_e$, dropping the months in $[A_e, E_e)$. The main
specification is
\begin{equation}
\ln y_{imt}^{e} = \beta\, T_i^{e} \times \text{Post}_t^{e} + \gamma\, B_i^{e} \times \text{Post}_t^{e}
 + \delta\, X_{c(i)t} + \alpha_{i m}^{e} + \lambda_{t}^{e} + \varepsilon_{imt}^{e},
\label{eq:did}
\end{equation}
where $y$ is seats, flights, seats per flight, CO$_2$, CO$_2$ per seat or CO$_2$ per seat-kilometre at
airport $i$ in calendar month $m$ of month $t$; $T_i^e$ marks airports in the taxing country and
$B_i^e$ airports in the 50--300 km border group; $\text{Post}_t^e$ equals one from $E_e$ onward;
$X_{c(i)t}$ are log GDP and log population of the airport's country; $\alpha_{im}^e$ are airport-by-calendar-month-by-event fixed effects, which absorb each airport's seasonal profile within each stack; and
$\lambda_t^e$ are month-by-event fixed effects. Stacking events with their own fixed effects avoids the
negative weights of two-way fixed-effects estimators under staggered adoption \citep{cengiz2019,goodmanbacon2021}.
Standard errors are clustered by country, the level at which the tax varies. We report equal-weighted
estimates, which describe the typical airport, and seat-weighted estimates, which describe the national
total, and we interact the treatment with the within-country size tercile.

The buffer-and-ring treatment of foreign airports is the analogue of the 2--20 km buffer in
\citet{li2019subway}. Airports just across the border are the ones most likely to gain passengers who
drive to avoid the tax; treating them as controls would bias the estimate toward larger losses, and
dropping them would discard the leakage margin. Entering them as their own group identifies the leakage
coefficient $\gamma$ directly and keeps the control group clean \citep{butts2023}. The Dutch episode,
in which the tax was introduced in July 2008 and set to zero in July 2009, is analysed separately as an
introduction followed by a removal and excluded from the pooled stacks.

\subsection{Event time and sensitivity to pre-trends}

The identifying assumption is that, absent the tax, seats at taxed airports would have followed the
control path. We test it with event-time versions of \eqref{eq:did} in twelve-month bins so that
seasonality cannot enter: two pre-announcement bins ($[A{-}36, A{-}25]$ and $[A{-}24, A{-}13]$),
the reference bin $[A{-}12, A{-}1]$, and post bins $[E, E{+}11]$ and $[E{+}12, E{+}23]$, pooled and event by
event. For annual outcomes (GACI and its components, and the decomposition of CO$_2$) the stacks run
from three years before to one year after the effective year with airport-by-event and year-by-event
fixed effects. We complement the pre-trend tests with the sensitivity analysis of \citet{rambachan2023},
which reports the set of post-period estimates consistent with post-treatment violations of parallel
trends up to a multiple $\bar M$ of the largest pre-period violation.

\subsection{Synthetic difference-in-differences for national totals}

The airport-level estimates describe the distribution of the effect across airports. To ask whether a
country's total flying and its national network position changed, we estimate the synthetic
difference-in-differences of \citet{arkhangelsky2021} event by event on the country-by-month log of
national seats and on the country-by-year seat-weighted mean of log GACI. Donors are the coded European
countries without a tax event in the window; unit weights match the treated country's pre-announcement
path and time weights match the pre-period to the post-period; inference uses placebo estimates that
treat each donor in turn.

\subsection{Decomposing the emissions effect and the network channel}

Annual departure CO$_2$ at an airport satisfies the identity
\begin{equation}
\ln \text{CO}_2 = \ln(\text{destinations}) + \ln(\text{flights per destination}) + \ln(\text{CO}_2 \text{ per flight}),
\label{eq:decomp}
\end{equation}
and CO$_2$ per flight further equals seats per flight times stage length times emissions per
seat-kilometre. Estimating \eqref{eq:did} on each term splits the emissions effect into a network
margin (routes dropped), a frequency margin (flights on surviving routes) and an aircraft-and-distance
margin, with coefficients that add up by construction. To relate the tax effect to connectivity, we
estimate the within-airport elasticity of CO$_2$ to GACI with airport and country-by-year fixed
effects, and apply the decomposition of \citet{gelbach2016} to split the tax effect into the part that
runs through the change in GACI and the part that does not. We treat this as a descriptive channel
analysis rather than a causal mediation estimate: an instrument for connectivity that is excludable from
the emissions equation and strong in the tax stacks is not available (Appendix), so we do not claim
that the network share identifies a structural indirect effect.

\subsection{Why not a country panel}

A natural alternative is a country-year panel of national emissions on a tax indicator with country and
year fixed effects. We report in the Appendix that such a regression is not robust: taxing countries are
concentrated in Northern and Central Europe, whose air markets grew more slowly after 2010 than the
Southern and Eastern European leisure markets that make up the never-taxed group, for reasons unrelated
to the tax, and the sign of the country-level coefficient depends on whether country trends are
included. The stacked design with announcement timing, a short window, airport-by-calendar-month fixed
effects and the buffer group is what removes this confounding, which is why we follow the local
comparison of \citet{li2019subway} rather than a cross-country regression.

% ---------------------------------------------------------------------------------------------
% Main text, part 2 (Sections 4-5). Numbers from tables_tax.tex / _res_tax_*.csv (2026-10-03).
% ---------------------------------------------------------------------------------------------

\section{Results}
\label{sec:results}

\subsection{Seats: the periphery loses, the hubs do not}

Table~\ref{tab:seats} reports the first-year effect of a ticket tax on log scheduled seats. With equal
airport weights, seats at taxed airports fall by 10.3 percent relative to control airports (column 1,
$p<0.05$). With seat weights the effect is 4.1 percent (column 2, $p<0.10$). The difference is the
distribution of the effect across the size terciles in columns 4 to 6: airports in the smallest tercile of
the taxing country lose 24.0 percent of their seats, the middle tercile 1.6 percent (not significant)
and the largest tercile 4.8 percent (not significant). The pre-trend coefficient for the second
pre-announcement year is $-0.031$ with a standard error of 0.021, and the pooled event study in
Table~\ref{tab:eventtime} shows pre-period coefficients of $-0.051$ and $-0.031$ (both insignificant)
against post-period coefficients of $-0.124$ in the first tax year and $-0.153$ in the second (both
$p<0.01$). Event by event, the first-year coefficients are $-0.18$ for Germany, $-0.12$ for Austria,
$-0.13$ for Norway and $-0.16$ for Sweden, all significant, against $-0.01$ for the UK rate increase and
$-0.10$ (not significant) for the Netherlands; the UK pre-period shows a trend and the UK event carries
the smallest dose.

The border group does not gain. Foreign airports 50 to 300 kilometres from the taxed country show
coefficients of $+0.009$ (equal weights) and $-0.030$ (seat weights), neither consistent with the
inflow of passengers that leakage would produce. Table~\ref{tab:robust} shows that the equal-weighted
estimate lies between $-0.11$ and $-0.15$ across effective-date windows of 12 to 36 months, without the
two recession-year events, with announcement donuts of two or three years and after removing
airport-specific pre-trends.

\subsection{Network position: concentration, not contraction}

Table~\ref{tab:network} turns to the annual network outcomes. The connectivity index of taxed airports
falls by 1.6 percent (equal weights, $p<0.05$), 2.3 percent (seat weights, $p<0.01$) and 2.6 percent at
hubs with more than one million pre-year seats. The components show where the index moves: taxed
airports lose 8.9 percent of their destinations, 18.6 percent of eigenvector centrality and 42.7
percent of betweenness centrality. By size, the index falls by 2.3 percent in the small and mid terciles
and by an insignificant 0.3 percent in the large tercile; destinations fall by 13.5 and 10.6 percent in
the small and mid terciles against 3.4 percent (not significant) at large airports; eigenvector
centrality falls by 32 percent at small airports and 13 percent at mid and large airports. The annual
event study (Table~\ref{tab:eventtime}, Panel B) is the cleanest evidence in the paper: the GACI
coefficients are $-0.001$ (SE 0.008) three years and $0.000$ (SE 0.003) two years before the effective
year, $-0.016$ in the effective year and $-0.018$ the year after (both $p<0.01$). Border airports show
no change in connectivity.

Two features of these estimates matter for interpretation. First, GACI is standardised within each
year, so the fall is relative: taxed airports slipped in the world ranking while the ranking of untaxed
airports improved. Second, the losses in degree and eigenvector centrality at small and mid-sized
airports, with hubs holding their degree, describe a network that is thinning at the edges and
concentrating at the centre. The betweenness losses are large in proportional terms because the
betweenness of small airports is close to zero and volatile; the point estimate for large airports is
insignificant.

\subsection{Emissions fall where flights disappear}

Table~\ref{tab:emissions} reports the same specifications for emissions. Departure CO$_2$ at taxed
airports falls by 10.4 percent with equal weights ($p<0.05$) and by 2.1 percent with seat weights (SE
0.021, not significant). By size, small airports lose 21.2 percent of their emissions ($p<0.01$), mid
airports 3.1 percent and large airports 6.4 percent (both not significant). Emissions per seat do not
fall: the seat-weighted coefficient is $+0.021$ ($p<0.05$), and emissions per seat-kilometre rise by
1.3 percent seat-weighted ($p<0.10$) and by 5.2 percent at small airports ($p<0.05$). Stage length is
unchanged everywhere. The tax therefore removes flying without making the flying that remains cleaner;
if anything, the loss of thin routes raises the average emissions intensity of the survivors.

Table~\ref{tab:mechanism} shows how. Panel A applies the exact decomposition \eqref{eq:decomp} to the
annual stacks. The total CO$_2$ effect of $-0.100$ splits into $-0.097$ from destinations, $-0.010$
from flights per destination and $+0.007$ from CO$_2$ per flight; for small airports the split is
$-0.161$, $-0.052$ and $+0.027$ of a total of $-0.186$. The entire emissions reduction is route exit.
Panel B confirms the margins at monthly frequency: small airports lose 23.7 percent of their flights,
keep aircraft size unchanged and are 4.4 percentage points more likely to record a month with no
scheduled departure at all; large airports lose 7.3 percent of flights ($p<0.10$) and raise seats per
flight by 2.5 percent ($p<0.05$). A flat per-passenger tax raises the fare on a thin, low-fare route by
a larger fraction than on a trunk route, so airlines cut the thin routes; at hubs, where transfer
passengers are exempt and the tax is a small share of the fare, they consolidate frequencies onto
larger aircraft.

\subsection{Connectivity and emissions}

Table~\ref{tab:gaci_co2} links the two sets of results in the manner of Table~9 in
\citet{li2019subway}, which translates network-density changes into pollution changes with the
estimated elasticity. Panel A estimates the within-airport elasticity of emissions to connectivity: a
one percent increase in GACI is associated with a 4.7 percent increase in CO$_2$ with airport and year
fixed effects and 3.6 percent with airport and country-by-year fixed effects, a 0.5 percent increase in
seats per flight, a 0.7 percent increase in stage length and a 0.5 percent fall in emissions per
seat-kilometre: more connected airports emit more in total and less per seat-kilometre. Multiplying
the tax effect on GACI from Table~\ref{tab:network} ($-0.016$ to $-0.023$) by this elasticity implies a
CO$_2$ change of 6 to 11 percent, in line with the direct equal-weighted estimate of $-10.4$ percent in
Table~\ref{tab:emissions}. Panel B makes the same point airport by airport with the
\citet{gelbach2016} decomposition: for small airports the total effect of $-0.186$ splits into $-0.069$
through the change in GACI and $-0.116$ holding GACI fixed; for mid airports $-0.073$ runs through
GACI and the direct part is zero; for large airports the direct part ($-0.052$) dominates a small
network component. Emissions intensity shows no network component anywhere. Because the elasticity
in Panel A is not instrumented within the tax stacks, we read Panel B as a consistent accounting of the
channel rather than as a structural mediation estimate (Appendix).

\subsection{National totals barely move}

The airport-level estimates say that taxes redistribute flying away from the periphery. Whether they
reduce national flying is a different question, and Table~\ref{tab:sdid} answers it with synthetic
difference-in-differences on national totals. For seats, the estimates are $-0.005$ for Germany (placebo
$p=1.00$), $-0.038$ for Austria ($p=0.83$), $-0.009$ for Sweden ($p=1.00$), $-0.016$ for the UK
($p=0.80$) and $-0.076$ for the Netherlands ($p=0.31$); only Ireland ($-0.159$, $p<0.01$), Malta
($-0.207$, $p<0.01$) and Norway ($-0.155$, $p=0.07$) show national losses, the first two coinciding with
the deepest national recessions in the sample and the third with a high flat tax in a market dominated
by domestic flying. On the seat-weighted national connectivity index, no event moves more than 5
percent and none is significant at 5 percent. Taken with the airport results, the picture is that of a
tax that reshuffles seats and network position within the country toward its hubs while leaving the
national total, and the country's position in the world network, largely where they were.

\subsection{Leakage}

Table~\ref{tab:leakage} splits the border group by distance to the nearest treated airport. Foreign
airports 50 to 150 kilometres away show $-0.019$ (SE 0.042) and those 150 to 300 kilometres away
$-0.002$ (SE 0.030); the 0 to 50 kilometre ring, which holds few airports and shows pre-trends, is
$-0.202$ ($p<0.10$). Splitting the 150 kilometre ring by size gives $0.000$ for foreign hubs with more
than five million seats and $-0.037$ for secondary airports. The one case with measurable leakage is
the Netherlands: while the 2008 tax was in force, seats at foreign airports 50 to 150 kilometres from
the border rose by 16.5 percent ($p<0.01$), and 14.9 percent in the eighteen months after the tax was
set to zero, consistent with the shift of Dutch passengers to Weeze, D\"usseldorf and Brussels
documented by \citet{gordijn2011}. The Dutch case combines a high rate, a dense foreign airport system
within an hour's drive and a one-year horizon; the other eight events show no detectable displacement
across borders.

\subsection{Cost per tonne}

Table~\ref{tab:abatement} turns the first-year estimates into an accounting of abatement, revenue and
connectivity lost, in the spirit of the cost-benefit section of \citet{li2019subway}. Applying the size-specific CO$_2$ coefficients to pre-year emissions, the 72 small taxed airports avoided 42 kt of CO$_2$
in the first tax year while generating EUR 20 million of revenue (at a load factor of 0.8), or EUR 484
of revenue per tonne avoided; the 67 mid airports avoided 53 kt at EUR 2,598 per tonne; and the 70
large airports, whose coefficient is not statistically different from zero, account for 3,836 kt and EUR
972 per tonne at the point estimate. Over all treated airports the revenue per tonne is EUR 988, an
order of magnitude above the EU allowance price of EUR 5 to 25 in the years of these events and five to
ten times the social cost of carbon of EUR 100 to 200 used in current appraisal. The abatement is also
bought with connectivity: small airports lose 0.021 index units and one destination per kilotonne
avoided, large airports 0.00007 index units and 0.04 destinations. The cost of the abatement a ticket
tax delivers is paid in the connectivity of peripheral regions.

\section{Discussion and conclusion}
\label{sec:conclusion}

Using nine European ticket-tax events, a monthly panel of more than 6,000 airports, flight-stage
emissions and a network-centrality index of airport position, this paper asks whether air passenger
taxes shrink the air network or concentrate it. The answer is the second. In the first tax year, taxed
airports lose 10 percent of their seats and emissions with equal weights but 4 and 2 percent with seat
weights; the smallest third of airports in the taxing country lose a quarter of their seats, a seventh of
their destinations and a third of their eigenvector centrality, while the largest third lose nothing and
upgauge. National seat totals and national network position are unchanged in synthetic
difference-in-differences for six of nine events. Emissions fall only through route exit; emissions per
seat-kilometre do not fall and rise at the hubs. Leakage across borders is absent except in the Dutch
episode.

Three implications follow. First, the flat per-passenger design with a transfer exemption is a tax on
thin, short, low-fare routes. It is therefore a tax on peripheral connectivity rather than on carbon:
the routes it removes are the ones on which the tax is a large share of the fare, not the ones with
the highest emissions. A levy proportional to fuel or distance, or an allowance price on emissions,
would target the intensity margin that the ticket tax leaves untouched, as the EU ETS results of
\citet{fageda2022pricing} suggest, at the cost of the hub exposure that \citet{fageda2025hub}
document. Second, the revenue a ticket tax raises per tonne avoided is far above any social cost of
carbon, which is consistent with the taxes' fiscal origins; if the objective is emissions, the
instrument is poorly matched to it. Third, because the losses fall on small and mid-sized airports and
the gains in position accrue to hubs, the tax raises the concentration of the national network. Our
companion work finds that hub concentration is associated with wider regional disparities, so the
distributional cost of the tax is likely to be regional as well as sectoral.

The estimates have limits. They are first-year effects on scheduled capacity and schedule-based
emissions; passengers, load factors and the longer-run adjustment of airline networks are not observed.
Identification rests on 33 country clusters and eight to nine events, and we therefore report event-by-event estimates, placebo inference for the national totals and sensitivity to parallel-trend violations rather than a single pooled standard error. The connectivity index is a relative measure, so the network result is a statement about position, not about the size of the world network. Finally, the
connectivity channel in Section~4.4 is an accounting decomposition; a credible instrument for airport
connectivity within short tax windows would allow the indirect effect to be identified rather than
described. Within these limits, the results give a clear answer to the question in the title: ticket
taxes cut carbon at the edges of the network and leave the centre intact, and the carbon they cut is
paid for in the connectivity of the periphery.

\clearpage
\section*{Tables}
% ---- Ticket-tax paper, lean table set (generated by 111_tax_tables_lean.py; do not edit by hand) ----

\begin{table}[!htbp]\centering
\footnotesize
\captionsetup{labelfont = bf, textfont = it, justification = centering}
\caption{National Air Passenger Taxes, 1998--2018}
\label{tab:events}
\resizebox{\linewidth}{!}{%
\begin{tabular}{lllll}
\toprule
Country & Tax & Announced & Effective & EUR per departing passenger \\
\midrule
Denmark & Passenger tax extended to domestic flights & 1997-05-06 & 1998-01-01 & DKK 75 ($\approx$ 10) \\
Malta & Departure tax doubled & 2004-11-24 & 2005-08-01 & Lm 10 increase ($\approx$ 23) \\
United Kingdom & Air Passenger Duty doubled & 2006-12-06 & 2007-02-01 & GBP 5 / 20 increase ($\approx$ 7 / 30) \\
Netherlands & Ticket tax (repealed July 2009) & 2007-02-07 & 2008-07-01 & 11.25 / 45 \\
Ireland & Air Travel Tax (repealed April 2014) & 2008-10-14 & 2009-03-30 & 2 / 10 \\
Germany & Luftverkehrsteuer & 2010-06-07 & 2011-01-01 & 8 / 25 / 45 \\
Austria & Flugabgabe & 2010-10-23 & 2011-04-01 & 8 / 20 / 35 \\
Norway & Flypassasjeravgift & 2015-12-14 & 2016-06-01 & NOK 80 ($\approx$ 9) \\
Sweden & Flygskatt & 2017-06-08 & 2018-04-01 & SEK 60 / 250 / 400 ($\approx$ 6 / 26 / 41) \\
\bottomrule
\end{tabular}%
}
\vspace{0.35em}
\begin{minipage}{0.97\linewidth}\footnotesize
\textit{Notes:} Announcement is the first public government decision (coalition agreement, cabinet or budget statement, parliamentary vote), read from primary sources; effective is the first taxed departure date. Rates are statutory amounts at introduction by distance band (or the increase, for the UK, Denmark and Malta), converted at event-month exchange rates. All nine taxes exempt connecting passengers on a single ticket, so hub transfer traffic is untaxed. Control airports are those of the other coded European countries (Austria, Belgium, Croatia, Denmark, Finland, France, Germany, Greece, Hungary, Iceland, Ireland, Luxembourg, Malta, Netherlands, Norway, Portugal, Spain, Sweden, Switzerland, United Kingdom) with no tax event of their own within three years before the announcement and two years after the effective date; Italy is excluded because its surcharge taxes transfer passengers. The Netherlands enters the pooled estimates of Table~\ref{tab:seats} only through the robustness windows (Appendix Table~\ref{tab:robust}) because its tax lasted twelve months. Sources: national legal gazettes and budget documents.
\end{minipage}
\end{table}

\begin{table}[!htbp]\centering
\footnotesize
\captionsetup{labelfont = bf, textfont = it, justification = centering}
\caption{Ticket Taxes and Scheduled Seats in the First Year}
\label{tab:seats}
\resizebox{\linewidth}{!}{%
\begin{tabular}{lcccccc}
\toprule
 & (1) & (2) & (3) & (4) & (5) & (6) \\
 & All airports & All, seat-weighted & Hubs only & Small & Mid & Large \\
\midrule
Taxed country $\times$ post & $-$0.103** & $-$0.041* & -- & $-$0.240*** & $-$0.016 & $-$0.048 \\
 & (0.046) & (0.021) &  & (0.066) & (0.060) & (0.051) \\
Foreign airports 50--300 km $\times$ post & 0.009 & $-$0.030* & -- & \multicolumn{3}{c}{0.010} \\
 & (0.027) & (0.016) & & \multicolumn{3}{c}{(0.027)} \\
\addlinespace
Pre-trend: taxed $\times$ [A$-$24, A$-$13] & $-$0.031 (0.021) & & & & & \\
\addlinespace
$N$ & 112{,}876 & 112{,}306 & & \multicolumn{3}{c}{112{,}876} \\
Countries & 33 & 31 & & \multicolumn{3}{c}{33} \\
\bottomrule
\end{tabular}%
}
\vspace{0.35em}
\begin{minipage}{0.97\linewidth}\footnotesize
\textit{Notes:} Dependent variable $\ln$(monthly departing seats). Stacked difference-in-differences over the eight tax increases of Table~\ref{tab:events} (Netherlands excluded): treated airports are those of the taxing country, controls the airports of the coded European countries without a tax event in the window; foreign airports within 50--300 km of a treated airport form their own group and those within 50 km are dropped. Pre-period: 36 months before the announcement; post-period: first twelve months from the effective date; months between announcement and effective date dropped. Fixed effects: airport-by-calendar-month-by-event and month-by-event. Column (2) weights by seats in the year before the announcement and so measures the effect on national totals. Size terciles are defined within each treated country by pre-announcement seats (median 30 thousand, 168 thousand and 2.0 million seats; the top tercile carries 94\% of treated seats). The pre-trend row is the coefficient on the second pre-announcement year relative to the first (full event-time coefficients in Appendix Table~\ref{tab:eventtime}). Standard errors clustered by country in parentheses. *, **, *** denote significance at the 10\%, 5\%, and 1\% levels.
\end{minipage}
\end{table}

\begin{table}[!htbp]\centering
\footnotesize
\captionsetup{labelfont = bf, textfont = it, justification = centering}
\caption{Ticket Taxes and the Network Position of Taxed Airports}
\label{tab:network}
\resizebox{\linewidth}{!}{%
\begin{tabular}{lcccccc}
\toprule
 & (1) & (2) & (3) & (4) & (5) & (6) \\
 & All airports & All, seat-weighted & Hubs only & Small & Mid & Large \\
\midrule
$\ln$ GACI (connectivity index) & $-$0.016** & $-$0.023*** & $-$0.026** & $-$0.023*** & $-$0.022*** & $-$0.003 \\
 & (0.006) & (0.007) & (0.010) & (0.005) & (0.003) & (0.015) \\
$\ln$ destinations & $-$0.089** & -- & -- & $-$0.135** & $-$0.106*** & $-$0.034 \\
 & (0.039) &  &  & (0.062) & (0.023) & (0.054) \\
$\ln$ eigenvector centrality & $-$0.186*** & $-$0.074* & $-$0.138 & $-$0.319*** & $-$0.131** & $-$0.131** \\
 & (0.046) & (0.040) & (0.082) & (0.082) & (0.049) & (0.053) \\
$\ln$ betweenness centrality & $-$0.427** & $-$0.088*** & $-$0.175 & $-$0.930** & $-$0.253** & $-$0.162 \\
 & (0.165) & (0.031) & (0.116) & (0.406) & (0.095) & (0.109) \\
\addlinespace
Foreign airports 50--300 km $\times$ post, $\ln$ GACI & $-$0.002 (0.005) & & & & & \\
Pre-trend, $\ln$ GACI: taxed $\times$ [$E-2$] & 0.000 (0.003) & & & $-$0.005** (0.002) & 0.007 (0.005) & $-$0.001 (0.004) \\
\addlinespace
$N$ (airport-years) & 9{,}615 & 13{,}569 & 3{,}308 & \multicolumn{3}{c}{9{,}615} \\
\bottomrule
\end{tabular}%
}
\vspace{0.35em}
\begin{minipage}{0.97\linewidth}\footnotesize
\textit{Notes:} Annual outcomes from the GACI panel: the Global Airport Connectivity Index (normalised each year, so coefficients are changes in relative network position), the number of destinations, eigenvector centrality and $\ln(1+10^{4}\times$normalised betweenness$)$. Columns (1) and (4)--(6): stacks of the three calendar years before the announcement and the first effective year, airport-by-event and year-by-event fixed effects, eight events. Columns (2)--(3): effective-date stacks of two years before and after, nine events; hubs are airports with at least one million seats in the pre-year. Size terciles as in Table~\ref{tab:seats}. The pre-trend row reports the coefficient two years before the effective year relative to the year before (full event-time coefficients in Appendix Table~\ref{tab:eventtime}). Standard errors clustered by country in parentheses. *, **, *** denote significance at the 10\%, 5\%, and 1\% levels.
\end{minipage}
\end{table}

\begin{table}[!htbp]\centering
\footnotesize
\captionsetup{labelfont = bf, textfont = it, justification = centering}
\caption{Ticket Taxes and Aviation CO$_2$ in the First Year}
\label{tab:emissions}
\resizebox{\linewidth}{!}{%
\begin{tabular}{lcccccc}
\toprule
 & (1) & (2) & (3) & (4) & (5) & (6) \\
 & All airports & All, seat-weighted & -- & Small & Mid & Large \\
\midrule
$\ln$ CO$_2$ (departing flights) & $-$0.104** & $-$0.021 & -- & $-$0.212*** & $-$0.031 & $-$0.064 \\
 & (0.047) & (0.021) &  & (0.056) & (0.066) & (0.051) \\
$\ln$ CO$_2$ per seat & 0.001 & 0.021** & -- & 0.024 & $-$0.005 & $-$0.015 \\
 & (0.014) & (0.008) &  & (0.026) & (0.011) & (0.015) \\
$\ln$ CO$_2$ per seat-kilometre & 0.008 & 0.013* & -- & 0.052** & $-$0.006 & $-$0.012 \\
 & (0.011) & (0.007) &  & (0.024) & (0.013) & (0.016) \\
\addlinespace
Foreign airports 50--300 km $\times$ post, $\ln$ CO$_2$ & 0.008 (0.034) & $-$0.023 (0.019) & & & & \\
\addlinespace
$N$ & 99{,}669 & 99{,}129 & & \multicolumn{3}{c}{99{,}669} \\
Countries & 33 & 31 & & \multicolumn{3}{c}{33} \\
\bottomrule
\end{tabular}%
}
\vspace{0.35em}
\begin{minipage}{0.97\linewidth}\footnotesize
\textit{Notes:} Monthly departing-flight CO$_2$ (LTO plus cruise) from the emissions panel; specification, sample, weights and size terciles as in Table~\ref{tab:seats}. CO$_2$ per seat-kilometre is the fuel intensity of the remaining flights. Column (3) is left empty because the hub definition of Table~\ref{tab:network} applies to the annual panel. Standard errors clustered by country in parentheses. *, **, *** denote significance at the 10\%, 5\%, and 1\% levels.
\end{minipage}
\end{table}

\begin{table}[!htbp]\centering
\footnotesize
\captionsetup{labelfont = bf, textfont = it, justification = centering}
\caption{How Emissions Fall: Route Exit, Not Fewer or Smaller Flights}
\label{tab:mechanism}
\resizebox{\linewidth}{!}{%
\begin{tabular}{lcccc}
\toprule
 & All treated & Small & Mid & Large \\
\midrule
\multicolumn{5}{l}{\textit{Panel A. Exact decomposition of annual CO$_2$: $\ln\text{CO}_2 = \ln\text{destinations} + \ln(\text{flights per destination}) + \ln(\text{CO}_2\text{ per flight})$}} \\
$\ln$ CO$_2$ (total) & $-$0.100** & $-$0.186** & $-$0.065 & $-$0.064 \\
 & (0.046) & (0.079) & (0.043) & (0.040) \\
\quad $\ln$ destinations (network margin) & $-$0.097** & $-$0.161** & $-$0.083*** & $-$0.051 \\
 & (0.039) & (0.060) & (0.025) & (0.054) \\
\quad $\ln$ flights per destination (frequency margin) & $-$0.010 & $-$0.052 & 0.025 & $-$0.011 \\
 & (0.044) & (0.110) & (0.030) & (0.039) \\
\quad $\ln$ CO$_2$ per flight (aircraft and distance margin) & 0.007 & 0.027 & $-$0.007 & $-$0.002 \\
 & (0.022) & (0.053) & (0.017) & (0.019) \\
\addlinespace
\multicolumn{5}{l}{\textit{Panel B. Monthly margins}} \\
$\ln$ flights & -- & $-$0.237*** & $-$0.052 & $-$0.073* \\
 & & (0.080) & (0.065) & (0.042) \\
$\ln$ seats per flight & -- & $-$0.002 & 0.036 & 0.025** \\
 & & (0.039) & (0.022) & (0.011) \\
Service loss (0/1, at-risk months) & -- & 0.044** & $-$0.006 & $-$0.003 \\
 & & (0.021) & (0.007) & (0.007) \\
\bottomrule
\end{tabular}%
}
\vspace{0.35em}
\begin{minipage}{0.97\linewidth}\footnotesize
\textit{Notes:} Panel A: the three components add up to total CO$_2$ by construction (annual stacks of Table~\ref{tab:seats}, three pre-years and the effective year; destinations from the GACI panel). Panel B: monthly stacks of Table~\ref{tab:seats}; service loss is a linear probability model on airport-months with scheduled departures twelve months earlier (months with no departure count as zero seats). Standard errors clustered by country. *, **, *** denote significance at the 10\%, 5\%, and 1\% levels.
\end{minipage}
\end{table}

\begin{table}[!htbp]\centering
\footnotesize
\captionsetup{labelfont = bf, textfont = it, justification = centering}
\caption{Connectivity and Emissions: More Connected Airports Emit More in Total but Less per Seat-Kilometre}
\label{tab:gaci_co2}
\resizebox{\linewidth}{!}{%
\begin{tabular}{lcccccc}
\toprule
\multicolumn{7}{l}{\textit{Panel A. Within-airport elasticities to $\ln$ GACI, all airports 1996--2019}} \\
 & $\ln$ CO$_2$ & $\ln$ destinations & $\ln$ seats per flight & $\ln$ stage length & $\ln$ CO$_2$ per seat-km & $\ln$ CO$_2$ per seat \\
\midrule
Airport and year FE & 4.722*** & 4.396*** & 0.488*** & 0.662*** & $-$0.532*** & 0.130** \\
 & (0.289) & (0.206) & (0.087) & (0.099) & (0.082) & (0.058) \\
Airport and country-by-year FE & 3.630*** & 4.178*** & 0.512*** & 0.762*** & $-$0.499*** & 0.263*** \\
 & (0.282) & (0.311) & (0.072) & (0.083) & (0.072) & (0.058) \\
$N$ & \multicolumn{6}{c}{79{,}613} \\
\addlinespace
\multicolumn{7}{l}{\textit{Panel B. Decomposition of the tax effect: direct effect holding GACI fixed, and the part running through GACI}} \\
 & \multicolumn{3}{c}{$\ln$ CO$_2$ (total)} & \multicolumn{3}{c}{$\ln$ CO$_2$ per seat-km} \\
\cmidrule(lr){2-4}\cmidrule(lr){5-7}
 & Total & Direct & Via GACI & Total & Direct & Via GACI \\
Small airports & $-$0.186** & $-$0.116 & $-$0.069 & 0.027 & 0.014 & 0.013 \\
Mid airports & $-$0.065 & 0.007 & $-$0.073 & 0.008 & $-$0.006 & 0.013 \\
Large airports & $-$0.064 & $-$0.052* & $-$0.012 & $-$0.002 & $-$0.004 & 0.002 \\
\bottomrule
\end{tabular}%
}
\vspace{0.35em}
\begin{minipage}{0.97\linewidth}\footnotesize
\textit{Notes:} Panel A: coefficients of $\ln$ GACI in airport-year regressions with the stated fixed effects; descriptive within-airport associations, not causal effects. Panel B: Gelbach (2016) decomposition on the annual tax stacks of Table~\ref{tab:mechanism}: the total tax coefficient by size tercile, the coefficient when $\ln$ GACI is added as a control (direct), and the difference, which equals the tax effect on $\ln$ GACI times the conditional coefficient of $\ln$ GACI (3.40 for total CO$_2$, $-$0.63 for CO$_2$ per seat-km). The decomposition is an accounting identity and assumes GACI is the only omitted channel; it is not an instrumental-variables mediation estimate. Standard errors clustered by country. *, **, *** denote significance at the 10\%, 5\%, and 1\% levels.
\end{minipage}
\end{table}

\begin{table}[!htbp]\centering
\footnotesize
\captionsetup{labelfont = bf, textfont = it, justification = centering}
\caption{Synthetic Difference-in-Differences by Event: National Totals Barely Move}
\label{tab:sdid}
\begin{threeparttable}
\begin{tabular}{lrcccc}
\toprule
 & & \multicolumn{2}{c}{National seats} & \multicolumn{2}{c}{Mean GACI (seat-weighted)} \\
\cmidrule(lr){3-4}\cmidrule(lr){5-6}
Event & Donors & SDID & Placebo $p$ & SDID & Placebo $p$ \\
\midrule
NLD 2008-07-01 & 13 & $-$0.076 & 0.31 & 0.020 & 0.54 \\
IRL 2009-03-30 & 13 & $-$0.159 & 0.00 & 0.006 & 0.92 \\
DEU 2011-01-01 & 11 & $-$0.005 & 1.00 & $-$0.009 & 0.82 \\
AUT 2011-04-01 & 12 & $-$0.038 & 0.83 & $-$0.019 & 0.58 \\
NOR 2016-06-01 & 14 & $-$0.155 & 0.07 & $-$0.019 & 0.38 \\
SWE 2018-04-01 & 13 & $-$0.009 & 1.00 & $-$0.047 & 0.08 \\
GBR 2007-02-01 & 15 & $-$0.016 & 0.80 & $-$0.033 & 0.29 \\
DNK 1998-01-01 & 16 & $-$0.099 & 0.25 & -- & -- \\
MLT 2005-08-01 & 13 & $-$0.207 & 0.00 & $-$0.053 & 0.15 \\
\bottomrule
\end{tabular}
\begin{tablenotes}\footnotesize
\item Notes: Country-level synthetic difference-in-differences (Arkhangelsky et al.\ 2021): unit and time weights on the simplex with the paper's ridge penalty. Seats: country-by-month $\ln$(seats) de-seasonalised with pre-announcement country-by-calendar-month means, 36 pre-months, first twelve effective months, anticipation months dropped. GACI: seat-weighted country mean of $\ln$ GACI, annual, six pre-years, effective year and the next (Denmark has too few pre-years). Donors: coded European countries without a tax event in the window. Placebo $p$: share of donors whose own estimate, treating that donor with the actual treated country removed, is at least as large in absolute value. Malta has two airports.
\end{tablenotes}
\end{threeparttable}
\end{table}

\begin{sidewaystable}[!htbp]\centering
\footnotesize
\captionsetup{labelfont = bf, textfont = it, justification = centering}
\caption{First-Year Accounting: Emissions Avoided, Revenue Raised and Connectivity Lost, by Airport Size}
\label{tab:abatement}
\resizebox{\linewidth}{!}{%
\begin{tabular}{lrrrrrr}
\toprule
 & Airports & CO$_2$ avoided (kt) & Revenue (m EUR) & EUR per t CO$_2$ & GACI lost (index units) & Destinations lost \\
\midrule
Small airports & 72 & 42 & 20 & 484 & 0.87 & 43 \\
Mid airports & 67 & 53 & 138 & 2{,}598 & 1.00 & 94 \\
Large airports & 70 & 3{,}836 & 3{,}728 & 972 & 0.26 & 168 \\
\midrule
All treated airports & 209 & 3{,}932 & 3{,}886 & 988 & 2.13 & 304 \\
\bottomrule
\end{tabular}%
}
\vspace{0.35em}
\begin{minipage}{0.97\linewidth}\footnotesize
\textit{Notes:} Each row applies the size tercile's first-year coefficients (CO$_2$ from Table~\ref{tab:emissions}; GACI and destinations from the annual stacks) to the pooled treated airports' pre-announcement-year CO$_2$, GACI and destinations. Revenue is the statutory tax at the airport's distance band times remaining passengers, assuming a load factor of 0.80. The large-airport coefficients are not statistically different from zero, so the large-airport and total rows are point estimates with wide uncertainty; the small- and mid-airport rows rest on precisely estimated coefficients. EUR per tonne compares revenue with avoided emissions and is not a welfare measure.
\end{minipage}
\end{sidewaystable}

\clearpage
\appendix
\setcounter{table}{0}
\renewcommand{\thetable}{A\arabic{table}}

\begin{table}[!htbp]\centering
\footnotesize
\captionsetup{labelfont = bf, textfont = it, justification = centering}
\caption{Event-Time Coefficients: No Pre-Trends before the Tax}
\label{tab:eventtime}
\begin{threeparttable}
\begin{tabular}{lcccc}
\toprule
 & \multicolumn{2}{c}{Before} & \multicolumn{2}{c}{After} \\
\cmidrule(lr){2-3}\cmidrule(lr){4-5}
\multicolumn{5}{l}{\textit{Panel A. Seats, monthly (bins of twelve months; reference: the year before the announcement)}} \\
 & [A$-$36, A$-$25] & [A$-$24, A$-$13] & [E, E$+$11] & [E$+$12, E$+$23] \\
All treated airports & $-$0.051 & $-$0.031 & $-$0.124*** & $-$0.153*** \\
 & (0.045) & (0.021) & (0.032) & (0.040) \\
\addlinespace
\multicolumn{5}{l}{\textit{Panel B. Connectivity (GACI), annual (reference: the year before the effective year)}} \\
 & $E-3$ & $E-2$ & $E$ & $E+1$ \\
All treated airports & $-$0.001 & 0.000 & $-$0.016*** & $-$0.018*** \\
 & (0.008) & (0.003) & (0.005) & (0.003) \\
Small & 0.004 & $-$0.005** & $-$0.021*** & $-$0.021*** \\
 & (0.005) & (0.002) & (0.005) & (0.006) \\
Mid & $-$0.006 & 0.007 & $-$0.021*** & $-$0.022** \\
 & (0.013) & (0.005) & (0.006) & (0.010) \\
Large & $-$0.001 & $-$0.001 & $-$0.004 & $-$0.014 \\
 & (0.012) & (0.004) & (0.012) & (0.013) \\
\bottomrule
\end{tabular}
\begin{tablenotes}\footnotesize
\item Notes: Specification of Tables~\ref{tab:seats} and \ref{tab:network} with event-time bins instead of a single post indicator. Panel A: $\ln$(seats), pooled eight events, bins of twelve calendar months so that seasonality cannot enter, reference bin the twelve months before the announcement (A); post bins count from the effective date (E). Panel B: $\ln$(GACI), annual, reference the year before the effective year. Border airports enter as a separate group (not shown). Standard errors clustered by country (33). *, **, *** denote significance at the 10\%, 5\%, and 1\% levels.
\end{tablenotes}
\end{threeparttable}
\end{table}

\begin{table}[!htbp]\centering
\footnotesize
\captionsetup{labelfont = bf, textfont = it, justification = centering}
\caption{Robustness of the First-Year Seat Effect (Equal Airport Weights)}
\label{tab:robust}
\begin{threeparttable}
\begin{tabular}{lcrr}
\toprule
Specification & Taxed $\times$ post & $N$ & Countries \\
\midrule
Effective-date window $\pm$12 months, nine events & $-$0.115*** (0.024) & 66{,}818 & 27 \\
Effective-date window $\pm$18 months, nine events & $-$0.130*** (0.030) & 93{,}915 & 27 \\
Effective-date window $\pm$24 months, nine events & $-$0.128*** (0.036) & 116{,}250 & 27 \\
Effective-date window $\pm$36 months, nine events & $-$0.130** (0.052) & 158{,}773 & 27 \\
Without the recession-year events (Netherlands 2008, Ireland 2009) & $-$0.111*** (0.025) & 51{,}894 & 27 \\
Announcement donut, two pre-years & $-$0.128*** (0.040) & 111{,}063 & 27 \\
Donut, airport-specific pre-trends removed & $-$0.153*** (0.036) & 107{,}904 & 27 \\
Donut, three pre-years, verified announcement dates (main) & $-$0.108** (0.050) & 95{,}270 & 33 \\
\bottomrule
\end{tabular}
\begin{tablenotes}\footnotesize
\item Notes: Treated-by-post coefficient of $\ln$(seats) from separate stacked regressions with the fixed effects of Table~\ref{tab:seats}. Window rows use the effective date with symmetric windows. Donut rows drop the months between announcement and effective date; the pre-trend row removes each airport's linear trend fitted on the 24 months before the announcement. The first four donut rows were estimated with the announcement dates as coded before verification (Austria 2011-01, Malta 2004-12); the main row uses the verified dates. Standard errors clustered by country. *, **, *** denote significance at the 10\%, 5\%, and 1\% levels.
\end{tablenotes}
\end{threeparttable}
\end{table}

\begin{table}[!htbp]\centering
\footnotesize
\captionsetup{labelfont = bf, textfont = it, justification = centering}
\caption{Cross-Border Leakage: Seats at Foreign Airports near the Taxed Country}
\label{tab:leakage}
\begin{threeparttable}
\begin{tabular}{lcc}
\toprule
 & (1) Distance rings & (2) Border airport type \\
\midrule
Taxed country $\times$ post & $-$0.108** (0.050) & $-$0.108** (0.050) \\
Foreign airports 0--50 km $\times$ post & $-$0.202* (0.100) &  \\
Foreign airports 50--150 km $\times$ post & $-$0.019 (0.042) &  \\
Foreign airports 150--300 km $\times$ post & $-$0.002 (0.030) & $-$0.002 (0.030) \\
Foreign hubs ($\geq$ 5 million seats) within 150 km $\times$ post & & 0.000 (0.025) \\
Foreign secondary airports within 150 km $\times$ post & & $-$0.037 (0.047) \\
\addlinespace
$N$ & 95{,}270 & 95{,}270 \\
\bottomrule
\end{tabular}
\begin{tablenotes}\footnotesize
\item Notes: Specification of Table~\ref{tab:seats}, column (1), with the border group split by great-circle distance to the nearest treated airport, or by size within 150 km. The 0--50 km ring holds few airports and shows pre-trends in the event study, so it is excluded elsewhere. The Dutch tax of 2008, analysed separately, is the one case with measurable leakage: seats at foreign airports within 150 km rose 16\% while the tax was in force (results available in the replication files). Standard errors clustered by country. *, **, *** denote significance at the 10\%, 5\%, and 1\% levels.
\end{tablenotes}
\end{threeparttable}
\end{table}

\clearpage
\appendix
\section{Supplementary results}
\label{app:extra}
\noindent (To be assembled from the result files.) A1: country-year regression of national emissions on a tax
indicator and its sensitivity to country trends (Section~3.5). A2: IV mediation with a geography-by-foreign-growth
instrument for connectivity, first-stage strength and Anderson--Rubin sets (Section~3.4). A3: HonestDiD sensitivity
sets for seats and CO$_2$ (Section~3.2). A4: Dutch episode, introduction and removal.
\bibliographystyle{plainnat}
\bibliography{refs_tax}
\end{document}
